\section{相关工作}

\paragraph{精确张量网络收缩。}
精确张量网络收缩的经典研究主要围绕收缩顺序、收缩树、树宽以及切片策略展开。
这一路线的核心结论很明确：真正决定复杂度的往往不是单个张量本身，而是中间状态的规模，以及由此诱导出的空间和时间开销\cite{markov2008simulating,dumitrescu2018benchmarking,jakes2019carving,gray2021hyperoptimized}。
在此基础上，index slicing 又进一步说明，当一次完整收缩难以承受时，可以把部分自由度切出来，转化为许多更小的子任务来并行执行或分批完成\cite{huang2021efficient}。
这些工作为本文提供了两个直接启发：一是困难首先来自中间分隔结构的膨胀；二是按块组织计算可以成为完整收缩过程中的重要系统手段。
不过，这一方向的主要目标仍然是精确收缩，而不是在完整输出仍需显式生成的前提下进行近似优化。

\paragraph{近似收缩与中间状态压缩。}
当精确收缩代价过高时，近似张量网络方法通常通过压缩中间对象来控制复杂度。
在规则拓扑上，TRG 等工作已经表明：面对不可承受的中间状态，能否做出有效截断，是近似收缩能否成立的关键\cite{levin2007trg}。
近年来，这一思路逐渐从规则晶格走向一般图结构。
Pan 等人提出了适用于任意连通结构的一般性近似收缩算法\cite{pan2020contracting}；Gray 和 Chan 将近似直接纳入收缩规划过程，使压缩与收缩顺序能够联合优化\cite{gray2024hyperoptimized}；Ma 等人则发展了面向一般图的 density-matrix 风格近似框架\cite{ma2024approximate}。
这些工作说明，一般图上的近似收缩已经可以系统地开展。
相比已有面向一般图的近似收缩方法，本文关注的是开放腿张量网络的完整收缩问题。由于输出部分需要被完整生成，算法不能改变其形式，因此更合理的做法是在保留输出部分的同时，对网络内部的接口状态进行近似处理。

\paragraph{随机化低秩近似与草图。}
本文的方法也借鉴了随机化数值线性代数与草图的思想。
随机化低秩近似表明，在许多大规模问题中，可以先用随机投影快速捕捉主导子空间，再在较小的子空间内做稳定分解，而不必一开始就构造昂贵的精确分解\cite{halko2011finding}。
更一般地，Count Sketch 与 Tensor Sketch 等方法展示了随机线性映射在高维对象压缩中的作用\cite{charikar2004finding,pham2013fast}。
非常近期的工作也开始直接讨论张量网络收缩中的草图近似\cite{heddes2026approximating}。
在本文中，随机化压缩是分隔状态框架中的实现工具之一，其作用是降低内部接口压缩的代价，而不是改变本文的问题设定。

\paragraph{与图模型消去方法的关系。}
从算法结构上看，本文也与图模型中的 bucket elimination 和 mini-bucket 有相通之处\cite{dechter1996bucket,dechter1997mini}。
这些方法同样强调：真正向上层传播的不是局部对象的全部细节，而是与外部仍然相关的接口信息；当接口过大时，就需要在中间层面做受控近似。
我们把这一思想放在开放腿 TNC 任务下重新组织，并把输出边界保持视为方法定义的一部分。
